\clearpage
\chapter*{Đính chính cho Bài trình bày số 149}
\addcontentsline{toc}{chapter}{Đính chính cho Bài trình bày số 149}
\setcounter{footnote}{0}
\sid{e149}
\markboth{Đính chính 149}{Đính chính 149}
\src{298}

\noindent\textbf{1. Định lý đối ngẫu cho các bó đại số nhất quán.}

\begin{center}
[Xê-mi-na Bourbaki, t. 9, 1956/57, số 149, 25 tr.]
\end{center}

\noindent\textbf{Ghi chú về trang 14.} --- Như tôi đã nêu trong bài giảng
của mình tại Đại hội Toán học Quốc tế năm 1958,\footnote{
Grothendieck (Alexander). --- The cohomology theory of abstract algebraic
varieties, \emph{Proceedings of the International Congress of Mathematicians}
[1958, Edinburgh], tr. 103--118. --- Cambridge, at the University Press, 1960.}
các vấn đề được nêu ra ở đây nay đã được giải quyết hoàn toàn.

Độc giả sẽ tìm thấy thêm thông tin về đối ngẫu của các bó nhất quán trong
\emph{loco citato}\footnotemark[1], tr. 112--115, trong EGA III, phần 2
(đang chuẩn bị), và trong khi chờ đợi, trong SGA 1962. Một cách trình bày
có hệ thống hơn sẽ được tìm thấy trong một chương sau của EGA
(chương IX theo đề cương dự kiến).

\par
